\documentclass[11pt]{article}
\usepackage{amsmath,amssymb,amsthm}
\usepackage[margin=1in]{geometry}
\newtheorem{theorem}{Theorem}
\newtheorem{lemma}[theorem]{Lemma}
\newcommand{\lean}[1]{}

\title{Odd Numbers, Briefly}
\author{An ordinary Lean user}
\date{}

\begin{document}
\maketitle

\section{Results}

\begin{theorem}\label{thm:odd-add}\lean{Plain.odd_add_odd}
The sum of two odd natural numbers is even.
\end{theorem}

\begin{theorem}\label{thm:sum-odd}
For every $n$, the sum of the first $n$ odd numbers is $n^2$.
\end{theorem}

\begin{theorem}\label{thm:square}
The square of an odd number is odd.
\end{theorem}

\begin{lemma}\label{lem:open}
% \lean{Plain.odd_add_odd} (a commented-out tag is no link)
Every odd number greater than one has an odd prime factor.
\end{lemma}

\end{document}
